\textit{Catatan edisi: solusi sumber memuat rumus turunan yang rusak dan berhenti sebelum menghitung kedua turunan berarah. Perhitungan berikut memperbaiki dan melengkapinya.}

\aufzaehlungdrei{Tuliskan

\mathdisp {\rho=\sqrt{x^2+y^2}} { }

dan

\mavergleichskettedisp
{\vergleichskette
{ A(x,y,z)
}
{ =} { \begin{pmatrix} x-xR\rho^{-1} \\ y-yR\rho^{-1} \\ z \end{pmatrix}
}
{ } {
}
{ } {
}
{ } {
}
}
{}{}{.}
Dengan demikian,

\mavergleichskettedisp
{\vergleichskette
{ h(x,y,z)
}
{ =} { x^2+y^2-2R\sqrt{x^2+y^2}+z^2+R^2
}
{ } {
}
{ } {
}
{ } {
}
}
{}{}{}
dan

\mavergleichskettedisp
{\vergleichskette
{ T
}
{ =} { h^{-1}(r^2)
}
{  }{
}
{    }{
}
{   }{
}
}
{}{}{.}
Gradien $h$ adalah $2A$. Pada $T$ berlaku $\lVert A\rVert=r$, sehingga suatu
\definitionsverweis {medan normal satuan}{}{}
adalah

\mavergleichskettedisp
{\vergleichskette
{ N(x,y,z)
}
{ =} { \frac{1}{r}\begin{pmatrix} x-xR\rho^{-1} \\ y-yR\rho^{-1} \\ z \end{pmatrix}
}
{ } {
}
{ } {
}
{ } {
}
}
{}{}{.}
}{Ambil

\mathdisp {P=(R-r,0,0)} { . }

Untuk menghitung turunan, kita dapat memakai perpanjangan gradien ternormalisasi $\widehat N=A/\lVert A\rVert$ pada suatu lingkungan $P$. Di titik ini,

\mavergleichskettedisp
{\vergleichskette
{ A(P)
}
{ =} { \begin{pmatrix}-r\\0\\0\end{pmatrix}
}
{ ,\quad }{ D_vA(P)
}
{ =} { \begin{pmatrix}0\\-\frac{r}{R-r}\\0\end{pmatrix}
}
{ } {
}
}
{}{}{,}
dan $\langle A(P),D_vA(P)\rangle=0$. Oleh karena itu,

\mavergleichskettedisp
{\vergleichskette
{ D_vN(P)
}
{ =} { D_v\widehat N(P)
}
{ =} { \frac{1}{r}D_vA(P)
}
{ =} { \begin{pmatrix}0\\-\frac{1}{R-r}\\0\end{pmatrix}
}
}
{}{}{.}
}{Demikian pula,

\mavergleichskettedisp
{\vergleichskette
{ D_wA(P)
}
{ =} { \begin{pmatrix}0\\0\\1\end{pmatrix}
}
{ ,\quad }{ \langle A(P),D_wA(P)\rangle
}
{ =} { 0
}
{ } {
}
}
{}{}{.}
Jadi,

\mavergleichskettedisp
{\vergleichskette
{ D_wN(P)
}
{ =} { D_w\widehat N(P)
}
{ =} { \frac{1}{r}D_wA(P)
}
{ =} { \begin{pmatrix}0\\0\\\frac{1}{r}\end{pmatrix}
}
}
{}{}{.}
}

 [[Kategorie:Latexseite]]
